\documentclass[10pt,a4paper]{amsart}
\usepackage[margin=19mm]{geometry}
\usepackage{amsmath,amssymb,mathtools}
\usepackage[scr=rsfs,cal=euler]{mathalfa}
\usepackage{xfrac}
\usepackage[unicode,hidelinks,bookmarks=false]{hyperref}
\newcommand{\ie}{i.e.\ }
\newcommand{\cf}{cf.\ }
\newcommand{\resp}{respectively\ }
\newcommand{\Subsection}{\subsection}
\providecommand{\nobreakdash}{\nobreak\hbox{-}}
\setlength{\emergencystretch}{2em}
\multlinegap0pt
\usepackage{bm}
\usepackage{cleveref}
\newcommand\N{\ensuremath{\mathbb{N}}}
\newcommand{\Psx}{\Psi_{\bm x}}
\newcommand\R{\ensuremath{\mathbb{R}}}
\newcommand{\falg}{$f\!$-algebra}
\newcommand{\one}{\mathbb{1}}
\theoremstyle{plain}
\newtheorem{thm}{Theorem}
\newtheorem{lem}[thm]{Lemma}
\newtheorem{prop}[thm]{Proposition}
\newtheorem{cor}[thm]{Corollary}
\theoremstyle{definition}
\newtheorem{defn}[thm]{Definition}
\newtheorem{example}[thm]{Example}
\theoremstyle{remark}
\newtheorem*{rem}{Remark}
\title[Lattice-ordered algebras admitting a polynomial growth continuous function calculus]{Lattice-ordered algebras admitting a polynomial growth continuous function calculus — Corollary 3.7}
\author{David Mu\~noz-Lahoz}
\date{Source: arXiv:2604.20294v2 (CC BY 4.0)}
\begin{document}
\maketitle

\noindent\emph{Source and licence.} The delivered work is the article shown in the title above, version arXiv:2604.20294v2 (CC BY 4.0), distributed under the Creative Commons Attribution 4.0 International licence, as recorded in the arXiv licence metadata for this exact version and shown on the abstract page saved with this item. The full licence notice, which carries the licence URL and the address of that abstract page, is the bundled file \texttt{licenses/arxiv-2604.20294-CC-BY-4.0.txt}.\par\smallskip

\noindent\textit{Editorial scope.} Counted unit: the article's Corollary 3.7 (source0.tex lines 1029--1054, source label \texttt{cor:conditions}); statement and proof are the article's own text line for line, with no line omitted, substituted or re-flowed. Required context retained uncounted: the article's f-algebra and polynomial-growth-calculus definitions and its Theorem 3.1 (source0.tex lines 414--417, 443--447 and 760--780), which the counted statement and proof invoke. Editorial formatting only: an AMS \texttt{amsart} preamble that restates the article's own macro definitions (\texttt{\textbackslash N}, \texttt{\textbackslash Psx}, \texttt{\textbackslash R}, \texttt{\textbackslash falg}, \texttt{\textbackslash one}) and the theorem environments the retained lines use, with \texttt{bm} loaded for the article's bold math and \texttt{cleveref} loaded so that the article's \texttt{\textbackslash cref} cross-reference resolves; the article ships no class or style file; the article's own macro definitions that appear after its \texttt{\textbackslash begin\{document\}} are restated in the wrapper preamble; sequential numbering of the retained environments in order of appearance, whereas the article prints the counted corollary as Corollary 3.7; equation numbering is not reproduced and every cross-reference inside the retained text resolves to the wrapper's numbering. No citation occurs inside any retained line, so the wrapper carries no bibliography. No asset-inclusion command, no drawing environment and no image asset occurs in this package.


\section*{Retained context: eq:proj\_one (source0.tex lines 414--417)}

    \begin{equation}\label{eq:proj_one}
            \Psx(\one)=1_X \text{ and }
\Psx(\pi _i)=x_i\text{ for }i=1,\ldots ,n.
    \end{equation}


\section*{Retained context: prop:admitPG (source0.tex lines 443--447)}

\begin{prop}\label{prop:admitPG}
    Let $X$ be a lattice-ordered algebra that admits a polynomial
    growth continuous function calculus without constants. Then $X$ is a commutative
    \falg.
\end{prop}


\section*{Retained context: thm:conditions (source0.tex lines 760--780)}

\begin{thm}\label{thm:conditions}
    Let $X$ be an Archimedean lattice-ordered algebra with identity $1_X$,
    and let $\bm x = (x_1,\ldots ,x_n)\in X^{n}$, for some $n \in \N$.
    The following are equivalent:
    \begin{enumerate}
    \item There is a lattice-algebra homomorphism $\Psx\colon
        PG_n\to X$ with $\Psx(\one)=1_X$ and $\Psx(\pi _i)=x_i$
        for $i=1,\ldots ,n$.
    \item Let $e=1_X\vee |x_1|\vee \cdots \vee |x_n|$. There exists an
        $f\!$-subalgebra $Y$ of $X$ such that $1_X,x_1,\ldots ,x_n \in
        Y$ and, for every $m \in \N$, the norm $\|{\cdot }\|_{e^{m}}$
        is complete on $Y\cap I_{e^{m}}$.
    \item There exist $f\ge 1_X\vee |x_1|\vee \cdots \vee |x_n|$ and an
        $f\!$-subalgebra $Y$ of $X$ such that $1_X,x_1,\ldots ,x_n \in
        Y$ and, for every $m \in \N$, the norm $\|{\cdot }\|_{f^{m}}$
        is complete on $Y\cap I_{f^{m}}$.
    \end{enumerate}
    When one and hence all three of these conditions are satisfied,
    the lattice-algebra homomorphism $\Psx$ satisfying the conditions
    in \textit{(i)} is unique.
\end{thm}


\section{The counted unit: Corollary 3.7 and its complete original proof}

\begin{cor}\label{cor:conditions}
    Let $X$ be an Archimedean lattice-ordered algebra with identity
    $1_X$. Then
    $X$ admits a polynomial growth continuous function calculus if and
    only if $X$ is an \falg\ and, for every $x_1,\ldots ,x_n \in X$,
    there exists $f\ge 1_X\vee |x_1|\vee \cdots
    \vee |x_n|$ such that, for every $m \in \N$, the norm
    $\|{\cdot }\|_{f^{m}}$ is complete on the closed sublattice
    generated by $x_1,\ldots ,x_n$ in $I_{f^{m}}$.

    In this case, the polynomial growth continuous function calculus
    is unique (in the sense that for each $n \in \N$ and $\bm x \in
    X^{n}$ there is only one lattice-algebra homomorphism $\Psx\colon
    PG_n\to X$ which satisfies \eqref{eq:proj_one}), and
    \begin{equation}\label{eq:compo}
        \Psi _{(\Psi _{\bm x}(f_1),\ldots ,\Psi _{\bm x}(f_m))}(g)=\Psi
        _{\bm x}(g\circ (f_1\times \cdots \times f_m))
    \end{equation}
    for each $m,n \in \N$, $\bm x \in X^{n}$, $f_1,\ldots ,f_m \in
    PG_n$ and $g \in PG_{m}$, where $f_1\times \cdots \times f_m\colon
    \R^{n}\to \R^{m}$ is the function defined by
    \[
        (f_1\times \cdots \times f_m)(t)=(f_1(t),\ldots ,f_m(t)),\quad
        t \in \R^{n}.
    \]
\end{cor}

\begin{proof}
    If $X$ admits a polynomial growth continuous function calculus,
    then $X$ is an \falg\ by \cref{prop:admitPG}. And if $x_1,\ldots ,x_n \in X$, then
    $e=1_X\vee|x_1|\vee \cdots \vee |x_n| $ has the stated properties
    by \cref{thm:conditions}. The converse also follows directly from \cref{thm:conditions}.

    Uniqueness was proved in \cref{thm:conditions}. To show
    \eqref{eq:compo}, let $f_1,\ldots ,f_m \in PG_n$ and $g \in PG_m$
    be as stated. First we need to check that $g\circ (f_1\times
    \cdots \times f_m) \in PG_n$. Suppose $|g|\le M (\one + |\pi _1| + \cdots +
    |\pi _m|)^{N}$ for some $M\in \R_+$ and $N \in \N$. Then it is immediate
    to check that
    \[
    |g\circ (f_1\times \cdots \times f_m)|\le M(\one + |f_1| + \cdots
    + |f_m|)^{N}.
    \]
    Since $f_1,\ldots ,f_m \in PG_n$, the right hand side of this
    equation belongs to $PG_n$. Since $PG_n$ is an ideal in
    $C(\R^{n})$, and $g\circ (f_1\times \cdots \times f_m)$ is
    continuous, it follows that $g\circ (f_1\times \cdots \times f_m)
    \in PG_n$. Hence, it makes sense to evaluate $\Psx$ at $g\circ (f_1\times
    \cdots \times f_m)$. This is true of every $g \in PG_m$.
    It is immediate to check that the map
    \[
    \begin{array}{cccc}
            & PG_m & \longrightarrow & X \\
            & g & \longmapsto & \Psx(g\circ (f_1\times \cdots \times
            f_m)) \\
    \end{array}
    \]
    is a lattice-algebra homomorphism that maps $\pi _i$ to $\Psx
    (f_i)$ for $i=1,\ldots ,m$. By uniqueness, it must be the map $\Psi
    _{(\Psx(f_1),\ldots ,\Psx(f_m))}$.
\end{proof}

\end{document}
